% !TeX program = pdflatex
\documentclass[11pt,a4paper]{article}
\usepackage{../../recursos/latex/estilo-notas}

\begin{document}
\cabecalhoaula{08}{Métricas para Classificação}{AME §7.4 e cap.~9; ISLP §4.4--4.5}

%==============================================================================
\section*{Objetivos da aula}
%==============================================================================
Ao final desta aula o aluno deve ser capaz de:
\begin{itemize}[leftmargin=1.4cm]
  \item explicar por que a \emph{acurácia} é insuficiente, sobretudo em dados
        desbalanceados;
  \item ler uma \textbf{matriz de confusão} e calcular sensibilidade,
        especificidade, precisão, VPN e $F_1$;
  \item formular \textbf{perdas assimétricas} e derivar o \emph{corte ótimo} de um
        classificador \emph{plug-in};
  \item construir e interpretar a \textbf{curva ROC} e a \textbf{AUC};
  \item escolher estratégias para \textbf{dados desbalanceados} (corte, pesos,
        re-amostragem) e entender que o problema real são os \emph{custos};
  \item distinguir \emph{classificar bem} de \emph{estimar bem
        probabilidades} (regras de pontuação próprias).
\end{itemize}

\begin{emsala}
Continuação direta da Aula~07: lá vimos que o classificador de Bayes minimiza a
probabilidade de erro --- mas a probabilidade de erro nem sempre é o que importa.
Pergunta de abertura: \emph{``um teste de câncer que diz `saudável' para todo mundo
acerta 99\% das vezes. Ele é bom?''}
\end{emsala}

%==============================================================================
\section{Por que a acurácia não basta}
%==============================================================================
O risco 0--1 (Aula~07) trata \emph{todo} erro como igual. Mas:
\begin{itemize}[leftmargin=1.4cm,nosep]
  \item \textbf{Custos assimétricos:} num \emph{screening}, deixar de detectar um
        doente (falso negativo) é muito pior que um alarme falso (falso positivo).
  \item \textbf{Dados desbalanceados:} com 10 doentes em 1000, o classificador
        trivial $g\equiv 0$ tem acurácia 99\% --- e sensibilidade zero. Como
        métodos buscam aproximar o classificador de Bayes, eles \emph{tendem} a
        esse classificador inútil quando a classe é rara.
\end{itemize}

\begin{atencao}
A raiz do problema \textbf{não é} o desbalanceamento em si, mas o fato de que os
\emph{custos} dos erros são diferentes --- de modo que o classificador de Bayes
(que minimiza só a probabilidade de erro) deixa de ser o classificador desejado.
Tenha isso em mente: as ``correções'' abaixo são formas de codificar custos.
\end{atencao}

%==============================================================================
\section{Matriz de confusão e métricas}
%==============================================================================
Para classificação binária ($Y\in\{0,1\}$, ``1'' $=$ positivo), cruzamos predito
$\times$ verdadeiro:
\begin{center}\small
\renewcommand{\arraystretch}{1.3}
\begin{tabular}{@{}ll cc@{}}
\toprule
 & & \multicolumn{2}{c}{\textbf{Verdadeiro}} \\
 & & $Y=0$ & $Y=1$ \\
\midrule
\multirow{2}{*}{\textbf{Predito}} & $\ghat=0$ & VN (verd.\ neg.) & FN (falso neg.) \\
                                  & $\ghat=1$ & FP (falso pos.) & VP (verd.\ pos.) \\
\bottomrule
\end{tabular}
\end{center}
A partir dela definem-se (cada uma estima uma probabilidade condicional
populacional):
\begin{align*}
  \text{Sensibilidade (\emph{recall}):}\ & S=\frac{\text{VP}}{\text{VP}+\text{FN}}
    && \approx \Prob(\ghat=1\mid Y=1),\\
  \text{Especificidade:}\ & E=\frac{\text{VN}}{\text{VN}+\text{FP}}
    && \approx \Prob(\ghat=0\mid Y=0),\\
  \text{Precisão (VPP):}\ & P=\frac{\text{VP}}{\text{VP}+\text{FP}}
    && \approx \Prob(Y=1\mid \ghat=1),\\
  \text{Valor preditivo negativo:}\ & \text{VPN}=\frac{\text{VN}}{\text{VN}+\text{FN}}
    && \approx \Prob(Y=0\mid \ghat=0).
\end{align*}
A \textbf{estatística $F_1$} combina sensibilidade e precisão pela média harmônica:
\begin{equation}\label{eq:f1}
  F_1 = \frac{2}{1/S + 1/P} = \frac{2\,P\,S}{P+S}.
\end{equation}

\begin{ideia}
\emph{Sensibilidade} olha ``dos positivos reais, quantos peguei?''; \emph{precisão}
olha ``dos que chamei de positivo, quantos acertei?''. O classificador trivial
$g\equiv 0$ tem $S=0$ e $E=1$: a especificidade alta esconde que ele falha em
\emph{todos} os doentes. Por isso olhe \emph{várias} métricas, não uma só.
\end{ideia}

\begin{atencao}
Todas essas quantidades devem ser calculadas no conjunto de \emph{validação/teste},
nunca no treino (senão superestimam o desempenho --- Aula~03).
\end{atencao}

%==============================================================================
\section{Perdas assimétricas e o corte ótimo}
%==============================================================================
Para codificar custos diferentes, redefina o risco com pesos $l_0,l_1>0$:
\begin{equation}\label{eq:riscoassim}
  \risco(g) = l_1\,\Prob\big(Y\ne g(\X),\,Y=0\big)
            + l_0\,\Prob\big(Y\ne g(\X),\,Y=1\big),
\end{equation}
onde $l_0$ penaliza errar um positivo (FN) e $l_1$ penaliza errar um negativo (FP).

\begin{proposicao}[Corte ótimo sob perda assimétrica]\label{prop:corte}
O classificador que minimiza \eqref{eq:riscoassim} é
\[
  g(\x) = \1{\ \Prob(Y=1\mid\x) > \tfrac{l_1}{l_0+l_1}\ }.
\]
\end{proposicao}
\begin{proof}
Condicionando em $\x$, o custo esperado de prever $1$ é $l_1\Prob(Y=0\mid\x)$ e o de
prever $0$ é $l_0\Prob(Y=1\mid\x)$. Prevemos $1$ quando
$l_1\Prob(Y=0\mid\x) < l_0\Prob(Y=1\mid\x)$, isto é,
$l_1(1-\eta) < l_0\,\eta$ com $\eta=\Prob(Y=1\mid\x)$, o que dá
$\eta > l_1/(l_0+l_1)$. Como cada $\x$ é otimizado pontualmente, segue o resultado.
\end{proof}

\begin{ideia}
Eis a mensagem operacional: \textbf{mudar custos $=$ mudar o corte}. Com $l_0=l_1$
recuperamos o corte $1/2$ (Bayes). Se errar um positivo é mais caro ($l_0>l_1$), o
corte cai abaixo de $1/2$ --- classificamos como positivo com menos evidência. Um
caso notável: tomar $l_c=\pi_c$ leva ao corte $\widehat\Prob(Y=1)$ (a prevalência),
em vez de $1/2$.
\end{ideia}

%==============================================================================
\section{Escolhendo o corte: curva ROC e AUC}
%==============================================================================
Um classificador \emph{plug-in} probabilístico gera uma \emph{família} de
classificadores variando o corte $K$: $g_K(\x)=\1{\widehat\Prob(Y=1\mid\x)\ge K}$.
Cada $K$ dá um par (sensibilidade, especificidade).

A \textbf{curva ROC} (\emph{Receiver Operating Characteristic}) traça a
sensibilidade ($\Prob(\ghat=1\mid Y=1)$) contra $1-$especificidade (a taxa de falso
positivo) à medida que $K$ varia de $1$ a $0$. Propriedades:
\begin{itemize}[leftmargin=1.4cm,nosep]
  \item o canto superior esquerdo $(0,1)$ é o classificador perfeito;
  \item a diagonal é o ``palpite aleatório'';
  \item a \textbf{AUC} (área sob a curva) resume o desempenho em um número:
        $\text{AUC}=1$ é perfeito, $0{,}5$ é aleatório. Interpretação:
        probabilidade de o modelo dar \emph{score} maior a um positivo sorteado do
        que a um negativo sorteado.
\end{itemize}

\begin{ideia}
A ROC/AUC avalia o \emph{ordenamento} dado pelo modelo \textbf{independentemente do
corte} --- útil para comparar classificadores antes de fixar custos. Já escolhido o
contexto, fixa-se $K$: critério comum é maximizar
$\text{Sensibilidade}+\text{Especificidade}$ (ponto mais próximo do canto
$(0,1)$), ou usar o corte da Prop.~\ref{prop:corte} se os custos são conhecidos.
\end{ideia}

\begin{atencao}
A curva ROC também deve ser construída na \emph{validação/teste}. E cuidado: em
dados muito desbalanceados, a ROC pode parecer otimista --- a \emph{curva
precisão--recall} costuma ser mais informativa nesses casos.
\end{atencao}

%==============================================================================
\section{Lidando com dados desbalanceados}
%==============================================================================
Três famílias de estratégias (todas, no fundo, codificam custos):
\begin{description}[leftmargin=2.6cm,style=nextline]
  \item[Ajustar o corte] para métodos probabilísticos, basta mover $K$
        (Prop.~\ref{prop:corte}). Em geral é a solução mais limpa.
  \item[Reponderar observações] dar peso maior à classe rara no treino, de modo que
        sua contribuição na função objetivo aumente (\texttt{class\_weight} no
        \texttt{scikit-learn}). Útil quando o método não estima $\Prob(Y=1\mid\x)$.
  \item[Re-amostrar] alterar artificialmente a prevalência: \emph{oversampling} da
        minoria (ex.: SMOTE) ou \emph{undersampling} da maioria.
\end{description}

\begin{atencao}
Repetindo a mensagem central: o ``problema'' não é o desbalanceamento, são os
\emph{custos desiguais}. Se os custos forem realmente iguais, um modelo bem
calibrado com corte $1/2$ pode ser o correto, mesmo desbalanceado. Não aplique
re-amostragem reflexivamente.
\end{atencao}

%==============================================================================
\section{Classificar bem $\neq$ estimar bem probabilidades}
%==============================================================================
Um classificador \emph{plug-in} pode estar muito perto do classificador de Bayes
\textbf{sem} que $\widehat\Prob\approx\Prob$: para a decisão, basta que
$\widehat\Prob(Y=1\mid\x)-K$ tenha o \emph{mesmo sinal} que
$\Prob(Y=1\mid\x)-K$. Ou seja, boa classificação não garante boa estimativa de
probabilidade.

Quando precisamos das \emph{probabilidades} em si (para ordenar riscos, decidir com
custos variáveis, ou inferência conformal), avaliamos com \textbf{regras de
pontuação próprias}, que são minimizadas pela probabilidade verdadeira:
\begin{itemize}[leftmargin=1.4cm,nosep]
  \item \textbf{Entropia cruzada / \emph{log-loss}}:
        $-\frac1n\sum_i\big[Y_i\log\widehat p_i + (1-Y_i)\log(1-\widehat p_i)\big]$
        (é a log-verossimilhança negativa da logística);
  \item \textbf{Score de Brier}:
        $\frac1n\sum_i (\widehat p_i - Y_i)^2$ (erro quadrático das
        probabilidades).
\end{itemize}
Esses critérios são os indicados para escolher \emph{tuning parameters} de
classificadores probabilísticos (ex.: $\lambda$ da logística penalizada) quando o
objetivo é estimar bem $\Prob$.

\begin{ideia}
\textbf{Calibração}: um modelo é bem calibrado se, dentre os casos a que ele atribui
$\widehat p=0{,}7$, cerca de 70\% são de fato positivos. \emph{Log-loss} e Brier
penalizam a má calibração; acurácia e AUC, não. Escolha a métrica pelo objetivo
real do problema.
\end{ideia}

%==============================================================================
\section{Prática em Python}
%==============================================================================
\begin{lstlisting}
from sklearn.metrics import (confusion_matrix, classification_report,
                             roc_auc_score, RocCurveDisplay,
                             log_loss, brier_score_loss)

p = modelo.predict_proba(X_te)[:, 1]     # probabilidades estimadas
y_hat = (p >= 0.5).astype(int)           # corte; ajuste conforme os custos

print(confusion_matrix(y_te, y_hat))
print(classification_report(y_te, y_hat))   # precisao, recall, F1 por classe
print("AUC:", roc_auc_score(y_te, p))
print("log-loss:", log_loss(y_te, p),
      " Brier:", brier_score_loss(y_te, p))
RocCurveDisplay.from_predictions(y_te, p)

# desbalanceamento: reponderar a classe rara
# LogisticRegression(class_weight="balanced")
\end{lstlisting}
O notebook \texttt{Aula prática 08} e o material \texttt{MatConf.pdf} exploram a
matriz de confusão e a ROC no \texttt{bank\_train\_redux.csv}.

%==============================================================================
\section*{Resumo}
%==============================================================================
\begin{itemize}[leftmargin=1.4cm]
  \item Acurácia engana sob custos assimétricos/desbalanceamento; olhe a
        \textbf{matriz de confusão}.
  \item Sensibilidade, especificidade, precisão, VPN e $F_1$ \eqref{eq:f1} contam
        histórias diferentes.
  \item \textbf{Perda assimétrica} \eqref{eq:riscoassim} $\Rightarrow$ corte ótimo
        $l_1/(l_0+l_1)$ (Prop.~\ref{prop:corte}): mudar custos é mudar o corte.
  \item \textbf{ROC/AUC} avaliam o ordenamento independentemente do corte; fixe $K$
        por custos ou por sensibilidade$+$especificidade.
  \item Desbalanceamento: corte, pesos ou re-amostragem --- mas a raiz são os
        custos.
  \item Classificar bem $\neq$ estimar $\Prob$ bem: use \emph{log-loss}/Brier e
        pense em \textbf{calibração}.
\end{itemize}

\paragraph{Leitura recomendada.} [AME] §7.4 (matriz de confusão e métricas), §9.1
(perdas assimétricas, cortes, dados desbalanceados), §9.2 (classificação vs.\
estimação de probabilidades); [ISLP] §4.4--4.5 (matriz de confusão, ROC).

\end{document}
